% Part: proof-theory
% Chapter: proof-search
% Section: completeness

\documentclass[../../../include/open-logic-section]{subfiles}

\begin{document}

\olfileid{pt}{ps}{cpl}

\olsection{完备性}

现在我们来证明这个证明搜索算法是无懈可击的，也就是说，如果它中止或者永不终止，那么起始相继式~$\Gamma \Sequent \Delta$ 就没有推导。为了证明这一点，我们定义一个结构~$\Struct M$，使得对所有 $!A \in \Gamma$ 都有 $\Sat{M}{!A}$，并对所有 $!A \in \Delta$ 都有 $\Sat/{M}{!A}$。换言之，$\Gamma$ 中的所有公式在~$\Struct M$ 中为真，而 $\Delta$ 中的所有公式在~$\Struct M$ 中为假，即 $\Gamma \Sequent \Delta$ 不是有效的。由于 \Log{G3c} 是可靠的，所以 $\Gamma \Sequent \Delta$ 没有推导。

我们利用一对公式 $\Theta, \Xi$ 和一个常项符号集~$C$ 来定义 $\Struct{M}$。我们从证明搜索的失败（中止或永不终止）运行中的一条\emph{失败分支}中得到 $\Theta$、$\Xi$ 和~$C$。一条失败分支是一个相继式序列 $\Pi_n \Sequent \Lambda_n$ 和常项符号集序列~$C_n$，满足 $\Pi_0 \Sequent \Lambda_0$ 就是末尾相继式 $\Gamma \Sequent \Delta$，并且 $\Pi_{n+1} \Sequent \Lambda_{n+1}$ 是第~$n+1$ 阶段生成的 $\Pi_n \Sequent \Lambda_n$ 的一个前提，且算法未能为该前提找到推导。显然，对每个~$n$ 这样的前提都必须存在，否则算法就已经找到推导了。我们令 $C_n$ 为与~$\Pi_n \Sequent \Lambda_n$ 相关联的常项符号集。

如果算法在一个只包含原子公式的顶端相继式处中止，则以该相继式为终点的分支就是一条失败分支。如果算法永不终止，它会不断生成越来越大的树。你可以将其想象成一棵不断生长的无穷树（即算法运行无限多步后你会得到的那棵树）。这将是一棵无穷树，但它是有限分支的（每个节点只有一到两个子节点——即紧接其上方的相继式）。根据 Kőnig（柯尼希）引理，每棵无穷的有限分支树都必有一条无穷分支。在搜索算法永远运行的情况下，这样一条无穷分支就是一条失败分支。

如果 $\Pi_n \Sequent \Lambda_n$ 和 $C_n$ 的序列是一条失败分支，令 $\Theta = \bigcup_n \Pi_n$，$\Xi = \bigcup_n \Lambda_n$（去掉指标，并去掉公式的多次出现），以及 $C = \bigcup_n C_n$。

现在我们可以定义~$\Struct{M}$了。

\begin{enumerate}
\item 论域 $\Domain{M}$ 是由~$C$ 中的常项符号和 $\Gamma \Sequent \Delta$ 中的函数符号（如果有的话）构成的不含变元的项的集合。
\item 若 $c \in \Domain{M}$，则 $\Assign{c}{M} = c$。
\item 若 $f$ 是 $\Gamma \Sequent \Delta$ 中的一个 $m$ 元函数符号，且 $t_1$,
\dots, $t_m \in \Domain{M}$，则 $\Assign{f}{M}(t_1, \dots, t_m) =
f(t_1, \dots, t_n)$。
\item 若 $R$ 是 $\Gamma \Sequent
\Delta$ 中的一个 $m$ 元谓词符号，则 $\Assign{R}{M} = \Setabs{\tuple{t_1, \dots, t_m} \in
\Domain{M}^n}{R(t_1, \dots, t_n) \in \Theta}$。
\end{enumerate}

$\Struct{M}$ 是一个\emph{项模型}，因为它的论域是一个项的集合，并且对常项符号和函数符号的解释保证了 $\Value{t}{M} = t$。我们利用~$\Theta$ 中的原子公式来定义 $\Assign{R}{M}$，从而确保 $\Sat{M}{R(t_1, \dots, t_n)}$ 当且仅当 $R(t_1, \dots, t_n)
\in \Theta$。

\begin{lem}\ollabel{lem:truth}
对所有 $!A \in \Theta \cup \Xi$：
\begin{enumerate}
  \item 若 $!A \in \Theta$，则 $\Sat{M}{!A}$。
  \item 若 $!A \in \Xi$，则 $\Sat/{M}{!A}$。
\end{enumerate}
\end{lem}

\begin{proof}
  对 $\depth{!A}$ 进行归纳。

首先假设 $!A$ 是原子的，即 $!A \ident \lfalse$ 或 $!A
\ident R(t_1, \dots, t_m)$。

由于 $\Pi_n \Sequent \Lambda_n$ 是一条失败分支，没有任何 $\Pi_n$ 会包含 $\lfalse$（否则 $\Pi_n \Sequent \Lambda_n$ 就会是一个公理）。根据~$\Struct{M}$ 的定义，$\Sat{M}{R(t_1, \dots, t_m)}$ 当且仅当 $R(t_1, \dots, t_m) \in
\Theta$。综合起来我们有：若 $!A \in \Theta$，则 $\Sat{M}{!A}$。

如果 $!A$ 是原子的且 $!A \in \Pi_n$，那么 $!A \in \Pi_{n+1}$；如果 $!A \in \Lambda_n$，那么 $!A \in \Lambda_{n+1}$。（这是因为证明搜索算法生成后继相继式的方式如此。）因此，如果 $!A \in \Theta \cap \Xi$，那么对某个~$n$，有 $!A \in \Pi_n \cap \Lambda_n$。这意味着 $\Pi_n \Sequent \Lambda_n$ 是一个公理，而失败分支不能包含任何公理。所以，根据构造，对原子的~$!A$ 有 $!A \notin \Theta \cap \Xi$，进而如果 $!A \in \Xi$ 则 $!A \notin \Theta$。根据~$\Struct{M}$ 的定义，这意味着如果 $!A \in \Xi$ 则 $\Sat/{M}{!A}$。

对于归纳步骤，假设 (1) 和 (2) 对所有深度小于~$!A$ 的公式成立。我们通过分情况讨论来证明 $!A$ 满足 (1) 和 (2)。

首先注意，当 $!A^i$ 出现在 $\Pi_n \Sequent \Lambda_n$ 中时，若 $i$ 不是 $\Pi_n \Sequent \Lambda_n$ 中的最小索引，则 $!A^i$ 也会出现在 $\Pi_{n+1} \Sequent \Lambda_{n+1}$ 中。由于在每一阶段添加的顶端相继式中最小索引会增加，我们最终会达到一个 $\Pi_n \Sequent \Lambda_n$，其中 $!A^i$ 出现且 $i$ 是最小索引。在第 $n+1$ 阶段，算法归约~$!A^i$。换句话说，如果 $!A \in \Theta$，则对某个 $n$，有 $\Pi_n = \Pi_n',
!A^i$ 且 $i$ 是最小索引。而如果 $!A \in \Xi$，则对某个 $n$，有 $\Lambda_n = \Lambda_n', !A^i$ 且 $i$ 是最小索引。

\begin{enumerate}
  \item $!A \in \Theta$ 且 $!A \ident !B \land !C$。则对某个 $n$，有 $\Pi_n = \Pi_n', (!B \land !C)^i$。$\Pi_{n+1} \Sequent
  \Lambda_{n+1}$ 是~\LeftR{\land} 的对应前提，即 $\Pi_{n+1} = \Pi_n', !B^k, !C^{k+1}$。因此，$!B, !C \in \Theta$。由归纳假设，$\Sat{M}{!B}$ 且 $\Sat{M}{!C}$。根据 $\Sat{M}{}$ 的定义，有 $\Sat{M}{!B \land !C}$。

  \item $!A \in \Xi$ 且 $!A \ident !B \land !C$。则对某个 $n$，有 $\Lambda_n = \Lambda_n', !B \land !C$。$\Pi_{n+1} \Sequent
  \Lambda_{n+1}$ 是以 $\Pi_n \Sequent \Lambda'_n, !B \land !C$ 为结论的~\RightR{\land} 的两个前提之一，即 $\Lambda_{n+1} = \Lambda_n', !B^k$ 或 $\Lambda_{n+1} = \Lambda_n',
  !C^k$。因此，要么 $!B \in \Xi$，要么 $!C \in \Xi$。由归纳假设，要么 $\Sat/{M}{!B}$，要么 $\Sat/{M}{!C}$。根据 $\Sat{M}{}$ 的定义，有 $\Sat/{M}{!B \land !C}$。

  \item $!A \ident \lnot !B$、$!A \ident !B \lor !C$ 和 $!A \ident !B \lif !C$ 的情况类似，留作练习。

  \item $!A \in \Theta$ 且 $!A \ident \lexists[x][!B(x)]$。则对某个 $n$，有 $\Pi_n = \Pi_n', \lexists[x][!B(x)]^i$。$\Pi_{n+1}
  \Sequent \Lambda_{n+1}$ 是~\LeftR{\lexists} 的对应前提，即对某个~$c$ 有 $\Pi_{n+1} = \Pi_n', !B(c)^k$。因此，$!B(c) \in \Theta$ 且~$c \in C$。由归纳假设，有 $\Sat{M}{!B(c)}$。根据 $\Sat{M}{}$ 的定义，有 $\Sat{M}{\lexists[x][!B(x)]}$。

  \item $!A \in \Xi$ 且 $!A \ident \lexists[x][!B(x)]$。则对某个 $n$，有 $\Lambda_n = \Lambda_n', \lexists[x][!B(x)]^i$。由于每次归约 $\lexists[x][!B(x)]$ 时，$\lexists[x][!B(x)]$ 都会带有一个新索引保留在前提相继式的右侧，因此，如果它出现在失败分支上，那么 $\lexists[x][!B(x)]$ 就会在右侧被无限次归约。对每个项 $t \in \Domain{M}$，在失败分支上它最终都会成为“仅由~$C_n$ 中的常项构成且在对 $\lexists[x][!B(x)]$ 进行右侧归约时第一个尚未被使用过的项”。因此，对所有 $t \in \Domain{M}$，在某个~$n$ 下 $!B(t) \in \Lambda_n$，从而 $!B(t) \in \Xi$。由归纳假设，对所有 $t \in \Domain{M}$ 有 $\Sat/{M}{!B(t)}$。根据 $\Sat{M}{}$ 的定义，有 $\Sat/{M}{\lexists[x][!B(x)]}$。

  \item $!A \ident \lforall[x][!B(x)]$ 的情况类似，留作练习。
\end{enumerate}
\end{proof}

\begin{cor}\ollabel{cor:complete} \Log{G3c} 的证明搜索算法是完备的：如果它中止或永不终止，则起始相继式~$\Gamma \Sequent \Delta$ 不是有效的，因而没有推导。
\end{cor}

\begin{proof}
  如果算法中止或永不终止，则存在一个失败分支，可以从中定义出 $\Struct{M}$。根据\olref{lem:truth}，对所有 $!A \in \Gamma$，有 $\Sat{M}{!A}$；对所有 $!A \in \Delta$，有 $\Sat/{M}{!A}$。（回想 $\Pi_0 = \Gamma$ 且 $\Lambda_0 = \Delta$，所以 $\Gamma \subseteq \Theta$ 且 $\Delta \subseteq \Xi$。）因此 $\Gamma \Sequent \Delta$ 不是有效的。由~\Log{G3c} 的可靠性，$\Gamma \Sequent \Delta$ 没有推导。
\end{proof}

\begin{cor}
\Log{G3c} 是完备的：如果 $\Gamma \Sequent \Delta$ 是有效的，则 $\Log{G3c} \Proves \Gamma \Sequent \Delta$。
\end{cor}

\begin{proof}
  我们证明其逆否命题。假设 $\Log{G3c} \Proves/ \Gamma \Sequent \Delta$。那么证明搜索算法必定中止或永不终止，因为找不到推导。根据\olref{cor:complete}，$\Gamma \Sequent \Delta$ 不是有效的。
\end{proof}

\end{document}